\newpage
\phantomsection
\label{prf:ec-reciprocal}

\begin{remark*}[Return]
\hyperref[thm:ec-reciprocal]{Return to Theorem}
\end{remark*}

\begin{theorem*}[Reciprocal of a Nonzero Equivalence Class]
For every $[(a,b)]\in\mathbb{Q}$ with $a\neq 0$,
\[
[(a,b)]^{-1} = [(b,a)].
\]
\end{theorem*}

\begin{proof}
\textbf{Professional Standard Proof.}~\\
Let $[(a,b)]\in\mathbb{Q}$ with $a\neq 0$. Then $a,b\in\mathbb{Z}$ with
$a\neq 0$ and $b\neq 0$. Since $a\neq 0$, the pair $(b,a)$ is an
admissible fraction-pair, so $[(b,a)]\in\mathbb{Q}$.

Compute the product using the definition of multiplication:
\[
[(a,b)]\cdot[(b,a)]
= [(ab,\,ba)]
= [(ab,\,ab)].
\]
Since $ab \neq 0$, we check $(ab,ab)\sim(1,1)$: indeed $ab\cdot 1 = 1\cdot ab$.
Therefore $[(ab,ab)] = [(1,1)] = 1$.

So $[(a,b)]\cdot[(b,a)] = 1$. By the uniqueness of multiplicative inverses
in a field, $[(a,b)]^{-1} = [(b,a)]$.
\end{proof}

\begin{proof}
\textbf{Detailed Learning Proof.}~\\

\textbf{Step 1.} Verify that $[(b,a)]$ is a well-formed element of $\mathbb{Q}$.

Since $(a,b)\in W$, we have $b\neq 0$. By hypothesis $a\neq 0$. Therefore
$(b,a)\in W$ (the denominator $a$ is nonzero), so $[(b,a)]\in\mathbb{Q}$.

\textbf{Step 2.} Recall the definition of multiplicative inverse.

The multiplicative inverse of a nonzero element $x\in\mathbb{Q}$ is the
unique element $x^{-1}\in\mathbb{Q}$ satisfying
\[
x\cdot x^{-1} = 1 = [(1,1)].
\]
Multiplicative inverses are unique in any field. We verify this for
$x = [(a,b)]$ and $x^{-1} = [(b,a)]$.

\textbf{Step 3.} Verify that $[(a,b)]\neq 0$.

By the zero class criterion, $[(a,b)] = 0$ iff $a = 0$. By hypothesis
$a\neq 0$, so $[(a,b)]\neq 0$, confirming that its inverse is defined.

\textbf{Step 4.} Compute the product $[(a,b)]\cdot[(b,a)]$.

By the definition of rational multiplication,
\[
[(a,b)]\cdot[(b,a)]
= [(a\cdot b,\; b\cdot a)]
= [(ab,\; ab)].
\]

\textbf{Step 5.} Show that $[(ab,ab)] = [(1,1)]$.

By the definition of $\sim$, $(ab,ab)\sim(1,1)$ iff
\[
(ab)\cdot 1 = ab\cdot 1,
\]
which is $ab = ab$, trivially true. Since $ab\neq 0$ (both $a$ and $b$
are nonzero), the pair $(ab,ab)$ is admissible, and thus
\[
[(ab,ab)] = [(1,1)] = 1.
\]

\textbf{Step 6.} Conclude that $[(a,b)]^{-1} = [(b,a)]$.

We have shown $[(a,b)]\cdot[(b,a)] = 1$. By uniqueness of multiplicative
inverses in $\mathbb{Q}$,
\[
[(a,b)]^{-1} = [(b,a)]. \qedhere
\]
\end{proof}

\begin{remark*}[Proof structure]
The proof is a direct verification. The candidate $[(b,a)]$ is the
``swapped'' fraction-pair. The product $[(a,b)]\cdot[(b,a)]$ evaluates
to $[(ab,ab)]$ by the multiplication formula, and $(ab,ab)\sim(1,1)$ since
the cross-product $ab\cdot 1 = 1\cdot ab$ is trivially true. Uniqueness of
multiplicative inverses then pins down $[(b,a)]$ as the inverse.

Note the logical dependency: this proof invokes uniqueness of multiplicative
inverses, which is a field-structure fact. This is why the reciprocal theorem
properly belongs \emph{after} $\mathbb{Q}$ has been established as a field,
not inside the construction section.
\end{remark*}

\begin{remark*}[Dependencies]~\\
\begin{itemize}
  \item \hyperref[def:rational-operations]{Definition~\ref*{def:rational-operations}}
  \item \hyperref[def:fraction-pair-equality]{Definition~\ref*{def:fraction-pair-equality}}
  \item \hyperref[thm:ec-multiplication]{Theorem~\ref*{thm:ec-multiplication}}
  \item \hyperref[def:rational-distinguished]{Definition~\ref*{def:rational-distinguished}}
        \textnormal{($1 = [(1,1)]$)}
  \item \hyperref[thm:ec-zero-numerator]{Theorem~\ref*{thm:ec-zero-numerator}}
        \textnormal{(to confirm $[(a,b)]\neq 0$)}
  \item \hyperref[thm:uniqueness-of-multiplicative-inverse]{Theorem~\ref*{thm:uniqueness-of-multiplicative-inverse}}
        \textnormal{(uniqueness of multiplicative inverse in $\mathbb{Q}$)}
\end{itemize}
\end{remark*}
